\section{System Overview}\label{sec:overview}

\begin{figure*}[t]
\centering
% Fig.: the components of GhostVault and what they talk to.
\begin{tikzpicture}[
  font=\footnotesize,
  node distance=4mm,
  box/.style={draw=black!65, fill=white, rounded corners=1.2pt, align=center, inner sep=2.6pt,
              minimum height=8.5mm, text width=#1},
  box/.default=27mm,
  core/.style={box=#1, fill=black!5},
  dark/.style={box=#1, fill=black!78, draw=black!78, text=white},
  region/.style={draw=black!35, rounded corners=3pt, inner sep=3.2mm},
  title/.style={font=\scriptsize\scshape, text=black!60, anchor=south west},
  arr/.style={-{Latex[length=1.5mm,width=1.2mm]}, draw=black!65, semithick},
  both/.style={{Latex[length=1.5mm,width=1.2mm]}-{Latex[length=1.5mm,width=1.2mm]}, draw=black!65, semithick},
  note/.style={font=\scriptsize\itshape, text=black!60, align=center},
]
% --- the user's side -------------------------------------------------------------------------
\node[box=25mm] (gui) {Tk window\\\texttt{gui/}};
\node[box=25mm, below=3mm of gui] (cli) {Command line\\\texttt{app.py}};

% --- the core --------------------------------------------------------------------------------
\node[core=40mm, right=10mm of gui, yshift=-5.75mm, minimum height=21mm] (vault)
  {\textbf{Vault} \texttt{vault.py}\\[1pt]
   ghost, restore, clean up, recover,\\ rename, previews, Reset\\[2pt]
   \scriptsize journal: one record per ghost\\
   \scriptsize placement \texttt{placement.py}\\
   \scriptsize recovery files, backup \texttt{backup.py}};

\node[box=19mm, below=8mm of vault.south west, anchor=north west] (crypto)
  {\texttt{crypto.py}\\\texttt{recovery.py}\\\scriptsize AES-GCM, HKDF,\\\scriptsize Argon2id};
\node[box=19mm, right=2mm of crypto] (files)
  {\texttt{files.py}\\\scriptsize the original:\\\scriptsize open, read,\\\scriptsize remove safely};
\node[box=19mm, left=2mm of crypto] (store)
  {\texttt{store.py}\\\texttt{fsutil.py}\\\scriptsize records, durable\\\scriptsize writes, lock};

% --- GitHub layer ----------------------------------------------------------------------------
\node[core=33mm, right=9mm of vault, minimum height=21mm] (gh)
  {\textbf{GitHub layer} \texttt{github.py}\\[2pt]
   \scriptsize ResilientApi: retries, rate\\\scriptsize limits, write pacing\\[1pt]
   \scriptsize HttpTransport: one HTTPS\\\scriptsize connection per thread,\\\scriptsize proxies, redirects};
\node[box=33mm, below=8mm of gh] (keys)
  {\texttt{credentials.py}\quad \texttt{tls.py}\\\scriptsize tokens \quad certificate checks};

% --- outside GhostVault ----------------------------------------------------------------------
\node[dark=31mm, right=10mm of gh, minimum height=21mm] (github)
  {\textbf{GitHub}\\[1pt]\scriptsize \texttt{api.github.com}\\[3pt]
   \scriptsize private repositories:\\\scriptsize blobs, trees, commits,\\\scriptsize one branch each};
\node[box=33mm, below=5mm of keys, fill=black!3] (os)
  {\scriptsize the operating system's\\\scriptsize keychain \quad trust store};

\node[box=40mm, below=7mm of crypto.south, anchor=north, xshift=-10.5mm, fill=black!3] (folder)
  {\textbf{Data folder} \scriptsize (owner-only)\\\scriptsize records with keys, encrypted\\\scriptsize previews, settings, log, lock};
\node[box=25mm, anchor=north, fill=black!3] (yours) at ([xshift=6mm]files.south |- folder.north)
  {\textbf{Your files}\\\scriptsize originals,\\\scriptsize restored files};

% --- arrows ----------------------------------------------------------------------------------
\draw[arr] (gui.east) -- ++(4.5mm,0) |- ([yshift=4mm]vault.west);
\draw[arr] (cli.east) -- ++(4.5mm,0) |- ([yshift=-4mm]vault.west);
\draw[arr] (vault.east) -- node[above, note] {requests} (gh.west);
\draw[both] (gh.east) -- node[above, note] {HTTPS} node[below, note] {JSON} (github.west);
\draw[arr] (gh.south) -- (keys.north);
\draw[both] (keys.south) -- (os.north);
\draw[arr] (vault.south -| crypto.north) -- (crypto.north);
\draw[arr] ([xshift=-12mm]vault.south) -- (store.north);
\draw[arr] ([xshift=12mm]vault.south) -- (files.north);
\draw[both] (store.south) -- (store.south |- folder.north);
\draw[both] (files.south) -- (files.south |- yours.north);

% --- regions ---------------------------------------------------------------------------------
\begin{scope}[on background layer]
  \node[region, fit=(gui)(cli)(vault)(gh)(keys)(os)(store)(files)(folder)(yours)] (pc) {};
\end{scope}
\node[title] at (pc.north west) {this computer};
\end{tikzpicture}

\caption{The parts of GhostVault. Everything inside the frame runs on the user's computer;
GhostVault has no server of its own. The core (\code{Vault} and the modules below it) never imports
the user interface. Only the GitHub layer talks to the network, and only to \code{api.github.com}.}
\label{fig:architecture}
\end{figure*}

\subsection{The parts}
GhostVault has three parts, shown in \figref{fig:architecture}.

\emph{The application} is one Python package, \code{ghostvault}. A Tk window is the normal way to
use it; a command line offers recovery operations and a self-check. It runs from the Python
package or as a standalone application with its own Python inside (\secref{sec:builds}).

\emph{The data folder} is a per-user folder that only the user can read. It holds one record per
ghost, the settings, an encrypted cache of image previews and a log. Its location follows each
system's convention (\secref{sec:folder}).

\emph{The repositories} are one or more private GitHub repositories, each with a token that can
write to it. GhostVault reaches them through GitHub's REST API over HTTPS and stores its objects in
one folder of each, \code{ghostvault/}. It has no server, account or service of its own: if
GhostVault's authors disappeared, every ghost would remain restorable with the application, the
data folder and the repositories.

\subsection{What a ghost is}
A ghost is two things: its record in the data folder, and its chunks on GitHub. The record
(\lstref{lst:record}) holds the file's key; for every chunk, the repository it went to, its
random path, its nonce and the ids that identify its ciphertext; the original's name, location,
size, SHA-256, permissions and timestamps; and the journal state of the operation in progress.
The chunks are encrypted, by default padded to one size, and named at random.

The record is the only map from a ghost to its chunks and the only copy of its key. The gallery
in the window is simply the set of readable records; there is no separate index that could
disagree with them. Losing a record without a recovery bundle (\secref{sec:recovery}) loses the
ghost, whatever GitHub still holds.

\subsection{The life of a file}
\figref{fig:lifecycle} shows the two operations step by step. Ghosting plans the upload, records
the plan, uploads and commits the chunks, verifies them on GitHub, and only then removes the
original. Restoring records where the file will go, downloads and checks every chunk, writes and
publishes the file, and only then removes the chunks from GitHub. The grey tags are the journal
states written to disk before the step they precede; they are what start-up recovery reads after
an interruption (\secref{sec:journal}).

\begin{figure}[t]
\centering
% Fig.: the steps of ghosting and restoring, and the journal state written before each step.
\begin{tikzpicture}[
  font=\scriptsize,
  stage/.style={draw=black!65, fill=white, rounded corners=1pt, align=center, text width=37mm,
               inner sep=2.2pt, minimum height=7mm,
               execute at begin node={\hyphenpenalty=10000\exhyphenpenalty=10000}},
  head/.style={stage, fill=black!78, draw=black!78, text=white, font=\scriptsize\bfseries, minimum height=5mm},
  arr/.style={-{Latex[length=1.3mm,width=1mm]}, draw=black!60},
  state/.style={fill=black!8, draw=black!30, rounded corners=4pt, inner sep=1.4pt, font=\scriptsize\ttfamily},
]
\matrix[row sep=5.6mm, column sep=5mm] (m) {
  \node[head] (g0) {ghost}; &
  \node[head] (r0) {restore}; \\
  \node[stage] (g1) {accept the original: a regular file with one name, opened without following links}; &
  \node[stage] (r1) {record the destination folder and a hidden temporary name}; \\
  \node[stage] (g2) {record the plan: a random key, each chunk's repository and random path}; &
  \node[stage] (r2) {download chunks in parallel; check each one, decrypt, write in order}; \\
  \node[stage] (g3) {read, encrypt and upload each chunk as a blob that no commit refers to yet}; &
  \node[stage] (r3) {check the whole file's SHA-256; restore its attributes; flush it to disk}; \\
  \node[stage] (g4) {record every blob id, then make one commit per repository}; &
  \node[stage] (r4) {publish it under its final name without replacing any file}; \\
  \node[stage] (g5) {verify: every blob id is on the branch; read back the first chunk and one other per repository}; &
  \node[stage] (r5) {delete the chunks, one commit per repository; confirm with a fresh listing}; \\[5mm]
  \node[stage] (g6) {move the original aside, check it there, then delete it}; &
  \node[stage] (r6) {forget the record}; \\
};
\foreach \a/\b in {g0/g1, g1/g2, g2/g3, g3/g4, g4/g5, g5/g6, r0/r1, r1/r2, r2/r3, r3/r4, r4/r5, r5/r6}
  \draw[arr] (\a) -- (\b);
% journal states, written durably before the step below them (drawn on the arrow)
\node[state] at ($(g2.south)!0.5!(g3.north)$) {uploading};
\node[state] at ($(g5.south)!0.3!(g6.north)$) {stored};
\node[state] at ($(g5.south)!0.72!(g6.north)$) {removing-original};
\node[state, below=4mm of g6] (ready) {ready};
\draw[arr] (g6) -- (ready);
\node[state] at ($(r1.south)!0.5!(r2.north)$) {restoring};
\node[state] at ($(r4.south)!0.5!(r5.north)$) {cleanup};
\end{tikzpicture}

\caption{Ghosting and restoring a file. Each grey tag is the journal state written durably before
the step below it. Nothing is deleted on either side until the other copy has been checked.}
\label{fig:lifecycle}
\end{figure}
